\section{Una cota positiva para las dos últimas capas}\label{sec:positiva}

Aplicamos primero la doble cuenta en una sola dirección. La cota determinista
del número de operaciones inversas convierte una cota inferior para la media de
reparaciones directas en la cota inferior para $f_h$ del Teorema~\ref{thm:positiva}.
El corte temprano proporciona las reparaciones bajo el proceso de referencia.
Queda comprobar qué proporción de la capa conserva esa cota bajo la ley uniforme.

Sea $h\in\{1,2\}$ y sea $\mathcal B_n\subseteq\Om_h(n)$ el evento de que algún cruce vacío tenga
menos de $a_0n/2$ reparaciones legales de un movimiento (incluido el filtro $G$).

\emph{Bajo el proceso.} Sea $\omega$ una trayectoria exitosa en la que vale~\eqref{eq:corte}.
Para todo cruce finalmente vacío, el Lema~\ref{lem:perdida} con $\delta=\delta_0$ y~\eqref{eq:constantes}
dan
\[
 A^{rc}_{\mathrm{final}}\ge a_0n-4\delta_0n-4\ge a_0n/2
\]
para $n$ grande, pues $a_0>8\delta_0$. La cota~\eqref{eq:corte} se obtuvo simultáneamente para
todas las raíces antes de seleccionar las que permanecen vacías, de modo que no hay sesgo de
supervivencia. Por el Lema~\ref{lem:corte},
\[
 \QQ_h(S\cap\{\text{la colocación final está en }\mathcal B_n\})\le e^{-cn}.
\]
No se necesita que el fallo del proceso completo sea raro.

\emph{Transferencia.} Por el Teorema~\ref{thm:comparacion}, $\QQ_h(S)\ge e^{-d_n}$, luego
$\nu'_h(\mathcal B_n)\le e^{-cn+d_n}\le e^{-cn/2}$ para $n$ grande. El
Lema~\ref{lem:binaria} y $\KL{\mu_h}{\nu'_h}\le d_n=o(n)$ dan
\[
 \mu_h(\mathcal B_n)\le\frac{d_n+\log2}{cn-d_n}\longrightarrow0 .
\]
Es una convergencia a cero, no una cola exponencial bajo $\mu_h$.

\emph{Doble cuenta.} Cada colocación de $\Om_h$ tiene exactamente $h^2$ cruces vacíos, y
$d_{h,1}$ suma las reparaciones sobre todos ellos; en $\mathcal B_n^c$, $d_{h,1}\ge h^2a_0n/2$. Por
el Lema~\ref{lem:doblecuenta},
\[
 f_h\ge\frac{\E_{\mu_h}d_{h,1}}{h^2(n-h)}\ge\frac{a_0n}{2(n-h)}\,\mu_h(\mathcal B_n^c),
 \qquad\text{luego}\qquad \liminf_{n\to\infty}f_h(n)\ge\frac{a_0}2 .
\]
La normalización por los $h^2$ cruces evita perder un factor $h^2$. \qed

\begin{observacion}
El resultado es asintótico y no asigna una constante útil a tamaños pequeños; no se sustituye
$a_0$ por la media numérica del perfil. La prueba no altera el proceso de referencia, no
condiciona sus transiciones durante el análisis y no usa ningún producto de vacancias bajo
$\mu_h$. Aquí el número de reinas colocadas tras el corte es $m_{\delta_0}\asymp n$.
\end{observacion}
